\documentclass[11pt]{article}
\usepackage[margin=1in]{geometry}
\usepackage[T1]{fontenc}
\usepackage{lmodern}
\usepackage{amsmath}
\usepackage{booktabs}
\usepackage{xcolor}
\usepackage[colorlinks=true,linkcolor=blue,citecolor=blue,urlcolor=blue]{hyperref}
\usepackage[nameinlink,noabbrev]{cleveref}

\title{Hyperlinks, URLs, and Colored Text Stress Case}
\author{Test Corpus}
\date{}

\begin{document}
\maketitle

\begin{abstract}
This fixture exercises hyperlinked text produced by \texttt{hyperref} with
\texttt{colorlinks} enabled, bare URL rendering, named external links, inline
colored spans created with \texttt{xcolor}, a highlighted span produced with
\verb|\colorbox|, and a footnote whose body contains an external hyperlink. It
also checks internal cross-references resolved through \texttt{cleveref} and
\verb|\nameref|. The prose is ordinary enough that a text-difference comparison
of the converted document remains informative.
\end{abstract}

\section{Introduction}
\label{sec:introduction}

Document conversion has to preserve both the visible anchor text of a link and
its underlying target. The opening section links to an external project page,
\href{https://example.org}{the project home page}, and also shows a bare
address that should be rendered verbatim and remain clickable:
\url{https://example.org/path/to/resource}. Color is used for emphasis rather
than for meaning, so the phrase \textcolor{red}{flagged for review} appears in
red and the phrase \textcolor{blue}{confirmed in the second pass} appears in
blue. A short highlighted span, \colorbox{yellow}{retained verbatim}, marks
text that an editor wanted to keep without alteration.

The introduction also includes a footnote whose body carries its own external
hyperlink, so that link handling can be checked inside the note
apparatus.\footnote{The reference implementation is archived at
\href{https://example.org/archive}{the archive entry}, which mirrors the
material discussed in the main text.} The remainder of the paper is organized as
follows. \Cref{sec:methods} describes the construction of the example, and
\cref{sec:results} reports a small summary table. The methods section is titled
``\nameref{sec:methods}'', and the results section is titled
``\nameref{sec:results}''.

\section{Methods}
\label{sec:methods}

The synthetic example is built from a fixed list of records, each tagged with a
status that determines its display color. Forward and backward references are
included so that the converter must resolve labels in both directions:
\cref{sec:introduction} introduces the link forms, while \cref{sec:results}
below presents the tabulated outcome. An additional external link points to a
specification document, \href{https://example.org/spec}{the format
specification}, which is cited again in the discussion. Bare addresses such as
\url{https://example.org/methods} are interspersed with named links so that both
rendering paths are exercised within a single paragraph.

\section{Results}
\label{sec:results}

\Cref{tab:status} summarizes the record statuses. One cell uses an inline color
command so that colored text inside a table can be distinguished from colored
text in ordinary prose. The accompanying anchor in the final column is a named
external link rather than a bare address.

\begin{table}[htbp]
\centering
\caption{Record statuses with one colored cell and a named external link.}
\label{tab:status}
\begin{tabular}{lll}
\toprule
Record & Status & Reference \\
\midrule
R-001 & \textcolor{red}{Flagged}  & \href{https://example.org/r001}{detail page} \\
R-002 & Confirmed                 & \href{https://example.org/r002}{detail page} \\
R-003 & Confirmed                 & \href{https://example.org/r003}{detail page} \\
R-004 & Pending                   & \url{https://example.org/r004} \\
\bottomrule
\end{tabular}
\end{table}

\section{Discussion}
\label{sec:discussion}

The example combines several link and color constructs that frequently require
special handling in Word output. Colored prose produced with \verb|\textcolor|,
the highlighted span produced with \verb|\colorbox|, the bare URL produced with
\verb|\url|, and the named links produced with \verb|\href| all appear in the
same short document, alongside a footnote that itself contains a hyperlink. The
internal references resolved by \verb|\cref|, \verb|\Cref|, and \verb|\nameref|
close the document by pointing back to \cref{sec:methods} and
\cref{tab:status}, while the format specification referenced earlier,
\href{https://example.org/spec}{the format specification}, demonstrates a link
target that appears more than once.

\end{document}
